% ══════════════════════════════════════════════════════════════════════════
%  Chapter 1: Prerequisites
% ══════════════════════════════════════════════════════════════════════════

\chapter{Prerequisites}\label{ch:prereqs}

We collect the facts from linear algebra and probability that are used
throughout this guide. The reader who is comfortable with inner product
spaces, the Kronecker product, Markov's inequality, and the central limit
theorem may safely skip to \Cref{ch:boolean-fourier}.

Standard references: \citet{axler2024linear} or \citet{horn2012matrix} for
linear algebra; \citet{durrett2019probability} or \citet{feller1971introduction}
for probability.

% ──────────────────────────────────────────────────────────────────────
\section{Inner Product Spaces}\label{sec:inner-product-spaces}

\begin{definition}[Inner product]\label{def:inner-product}
Let $V$ be a real vector space. An \emph{inner product}\index{inner product} on $V$ is a
function $\ip{\cdot}{\cdot} \colon V \times V \to \R$ satisfying, for all
$u, v, w \in V$ and $\alpha \in \R$:
\begin{enumerate}[label=(\roman*)]
  \item \emph{Symmetry:} $\ip{u}{v} = \ip{v}{u}$.
  \item \emph{Linearity in the first argument:}
        $\ip{\alpha u + v}{w} = \alpha\ip{u}{w} + \ip{v}{w}$.
  \item \emph{Positive definiteness:} $\ip{v}{v} \geq 0$, with equality
        if and only if $v = 0$.
\end{enumerate}
The pair $(V, \ip{\cdot}{\cdot})$ is called an \emph{inner product space}.
\end{definition}

The inner product induces a norm $\norm{v} = \sqrt{\ip{v}{v}}$.
The key structural fact is:

\begin{theorem}[Cauchy--Schwarz inequality]\label{thm:cauchy-schwarz}
For all $u, v$ in an inner product space,
$|\ip{u}{v}| \leq \norm{u}\,\norm{v}$,
with equality if and only if $u$ and $v$ are linearly dependent.%
\footnote{Proof: see \citet[Chapter~6]{axler2024linear} or
\citet[Section~5.1]{horn2012matrix}.}
\end{theorem}

\begin{definition}[Orthonormal basis]\label{def:onb}
A set $\{e_1, \ldots, e_n\}$ in an inner product space $V$ is
\emph{orthonormal} if $\ip{e_i}{e_j} = \delta_{ij}$ (the Kronecker delta).
If additionally $\operatorname{span}\{e_1,\ldots,e_n\} = V$, it is an
\emph{orthonormal basis} (ONB).
\end{definition}

The central fact about orthonormal bases is:

\begin{proposition}[Fourier expansion in a finite-dimensional inner product space]\label{prop:fourier-finite}
If $\{e_1, \ldots, e_N\}$ is an ONB for $(V, \ip{\cdot}{\cdot})$, then
every $v \in V$ has a unique expansion
\[
  v = \sum_{i=1}^{N} \ip{v}{e_i}\, e_i,
\]
and \emph{Parseval's identity} holds:
$\norm{v}^2 = \sum_{i=1}^{N} \ip{v}{e_i}^2$.
\end{proposition}

\begin{proof}
Write $v = \sum_i c_i e_i$. Taking the inner product with $e_j$ gives
$\ip{v}{e_j} = c_j$ by orthonormality. Then
$\norm{v}^2 = \ip{v}{v} = \sum_i c_i^2$.
\end{proof}

This is exactly the abstract template for the Boolean Fourier expansion
in \Cref{ch:boolean-fourier}: the vector space will be the space of all
real-valued functions on $\pmone^n$, the inner product will be the
expectation inner product, and the ONB will be the parity functions.


% ──────────────────────────────────────────────────────────────────────
\section{The Kronecker Product}\label{sec:kronecker}

\begin{definition}[Kronecker product]\label{def:kronecker}
Let $A$ be an $m \times n$ matrix and $B$ a $p \times q$ matrix. Their
\emph{Kronecker product}\index{Kronecker product} $A \otimes B$ is the $mp \times nq$ block matrix
\[
  A \otimes B =
  \begin{pmatrix}
    a_{11} B & a_{12} B & \cdots & a_{1n} B \\
    a_{21} B & a_{22} B & \cdots & a_{2n} B \\
    \vdots & \vdots & \ddots & \vdots \\
    a_{m1} B & a_{m2} B & \cdots & a_{mn} B
  \end{pmatrix}.
\]
\end{definition}

Key properties (all straightforward to verify; see
\citealp[Chapter~4]{horn2012matrix} for proofs):
\begin{enumerate}[label=(\roman*)]
  \item $(A \otimes B)(C \otimes D) = (AC) \otimes (BD)$, when the
        products $AC$ and $BD$ are defined.
  \item $(A \otimes B)^\top = A^\top \otimes B^\top$.
  \item If $A$ and $B$ are orthogonal ($A^\top A = I$ and $B^\top B = I$),
        so is $A \otimes B$. (More generally, if $A^\top A = \alpha I$
        and $B^\top B = \beta I$, then
        $(A \otimes B)^\top(A \otimes B) = \alpha\beta \cdot I$.)
  \item $\operatorname{tr}(A \otimes B) = \operatorname{tr}(A)\operatorname{tr}(B)$.
\end{enumerate}

Property (i) is the \emph{mixed-product rule}; it is the algebraic engine
behind the recursive structure of the Hadamard matrix (\Cref{ch:wht}).

\begin{example}[Mixed-product rule]\label{ex:mixed-product}
Let $A = \bigl(\begin{smallmatrix} 1 & 1 \\ 1 & -1 \end{smallmatrix}\bigr) = H_1$.
Then $(H_1 \otimes H_1)(H_1 \otimes H_1) = (H_1 H_1) \otimes (H_1 H_1)
= (2I_2) \otimes (2I_2) = 4I_4$. This verifies that $H_2^2 = 4I_4$
(which we will prove in general in \Cref{ch:wht}).
\end{example}

\begin{example}\label{ex:kronecker-2x2}
The Hadamard matrix $H_1 = \bigl(\begin{smallmatrix} 1 & 1 \\ 1 & -1
\end{smallmatrix}\bigr)$ satisfies
\[
  H_1 \otimes H_1 =
  \begin{pmatrix}
    1 & 1 & 1 & 1 \\
    1 & -1 & 1 & -1 \\
    1 & 1 & -1 & -1 \\
    1 & -1 & -1 & 1
  \end{pmatrix} = H_2.
\]
See \Cref{sec:hadamard-def} for the general construction.
\end{example}


% ──────────────────────────────────────────────────────────────────────
\section{Probability Essentials}\label{sec:prob-essentials}

Throughout this guide, probability spaces are finite (typically uniform
on $\pmone^n$), so measure-theoretic subtleties do not arise.

\begin{definition}[Expectation and variance]\label{def:expectation}
For a random variable $X$ on a finite probability space
$(\Omega, \mu)$ (where $\mu(\omega)$ is the probability of outcome $\omega$):
\[
  \E[X] = \sum_{\omega \in \Omega} \mu(\omega)\, X(\omega),
  \qquad
  \Var[X] = \E[(X - \E[X])^2] = \E[X^2] - (\E[X])^2.
\]
\end{definition}

\begin{theorem}[Markov's inequality]\label{thm:markov}\index{Markov's inequality}
If $X \geq 0$ and $a > 0$, then $\Prob[X \geq a] \leq \E[X] / a$.
\end{theorem}

\begin{proof}
$\E[X] = \sum_\omega \mu(\omega) X(\omega) \geq \sum_{\omega : X(\omega) \geq a} \mu(\omega) \cdot a = a \cdot \Prob[X \geq a]$.
\end{proof}

\begin{corollary}[Chebyshev's inequality]\label{cor:chebyshev}\index{Chebyshev's inequality}
$\Prob[|X - \E[X]| \geq t] \leq \Var[X] / t^2$.
\end{corollary}

\begin{proof}
Apply Markov to $(X - \E[X])^2$ with $a = t^2$.
\end{proof}

Markov's inequality is used directly in the thesis to prove the
degree-truncation approximation bound (Proposition~3.1 of the thesis).


% ──────────────────────────────────────────────────────────────────────
\section{The Central Limit Theorem and Berry--Ess\'een}\label{sec:clt}

\begin{theorem}[Central Limit Theorem]\label{thm:clt}\index{central limit theorem}
Let $X_1, X_2, \ldots$ be independent, identically distributed random
variables with $\E[X_i] = 0$ and $\Var[X_i] = \sigma^2 > 0$. Write
$S_n = X_1 + \cdots + X_n$ for the partial sum. Then
\[
  \frac{S_n}{\sigma\sqrt{n}}
  \xrightarrow{d} \mathcal{N}(0, 1)
  \quad \text{as } n \to \infty.
\]
Here $\xrightarrow{d}$ denotes \emph{convergence in distribution}:
$\Prob[S_n / (\sigma\sqrt{n}) \leq t] \to \Phi(t)$ for every $t$,
where $\Phi(t) = \frac{1}{\sqrt{2\pi}}\int_{-\infty}^t e^{-s^2/2}\,ds$
is the standard normal CDF.%
\footnote{For proof via characteristic functions, see
\citet[Chapter~3]{durrett2019probability} or
\citet[Chapter~XV]{feller1971introduction}.}
\end{theorem}

The CLT tells us that $S_n / (\sigma \sqrt{n})$ is approximately
$\mathcal{N}(0,1)$ for large $n$, but says nothing about the rate of
convergence. The Berry--Ess\'een theorem quantifies this.

\begin{theorem}[Berry--Ess\'een; \citealp{berry1941accuracy,esseen1942rate}]\label{thm:berry-esseen}\index{Berry--Ess\'een theorem}
Under the hypotheses of \Cref{thm:clt}, with additionally
$\E[|X_i|^3] = \gamma < \infty$, we have
\[
  \sup_{t \in \R}
  \left|
    \Prob\!\left[\frac{S_n}{\sigma\sqrt{n}} \leq t\right]
    - \Phi(t)
  \right|
  \leq \frac{C \gamma}{\sigma^3 \sqrt{n}},
\]
where $C$ is a universal constant (the best known value is $C < 0.4748$;
see \citealp{shevtsova2011absolute}).
\end{theorem}

\paragraph{Why this matters for the thesis.}
In the thesis, binary neural network layers compute sums of the form
$y_j = \sum_{i=1}^d w_{ji} x_i$ where $w_{ji}, x_i \in \pmone$.
When $x$ is uniform on $\pmone^d$, the summands $w_{ji} x_i$ are
independent $\pm 1$ random variables (regardless of the fixed signs
$w_{ji}$), so $\E[w_{ji} x_i] = 0$, $\Var[w_{ji} x_i] = 1$, and
$\E[|w_{ji} x_i|^3] = 1$. By Berry--Ess\'een,
$y_j / \sqrt{d}$ is within $O(1/\sqrt{d})$ of $\mathcal{N}(0,1)$ in CDF.
At $d = 1024$, this is about $0.015$---excellent approximation.
This justifies the $1/\sqrt{d}$ normalization used throughout the thesis
as a replacement for BitNet's learned scaling factors.


% ──────────────────────────────────────────────────────────────────────
\section{$L^p$ Norms}\label{sec:lp-norms}

For a function $f$ on a probability space $(\Omega, \mu)$:

\begin{definition}[$L^p$ norm]\label{def:lp}
For $p \geq 1$:
\[
  \norm{f}_p = \left(\E\!\left[|f|^p\right]\right)^{1/p}
  = \left(\sum_{\omega \in \Omega} \mu(\omega)\, |f(\omega)|^p\right)^{1/p}.
\]
For $p = \infty$: $\norm{f}_\infty = \max_\omega |f(\omega)|$.
\end{definition}

On a probability space, we always have $\norm{f}_p \leq \norm{f}_q$ for
$1 \leq p \leq q$. To see this: the function $t \mapsto t^{q/p}$ is
convex (since $q/p \geq 1$), so by Jensen's inequality
(see, e.g., \citealp[Theorem~1.6.2]{durrett2019probability}),
$(\E[|f|^p])^{q/p} \leq \E[|f|^q]$; taking $q$-th roots gives
$\norm{f}_p \leq \norm{f}_q$. (The direction seems backwards compared
to ``bigger exponent = bigger norm'' for sums, but the normalization by
the probability measure reverses things.)

The \emph{hypercontractivity} results in
\Cref{ch:concentration} are striking precisely because they give a reverse
inequality $\norm{T_\rho f}_q \leq \norm{f}_p$ for $q > p$, at the cost
of applying the noise operator $T_\rho$.

\begin{remark}[Notation]
Throughout this guide, when we write $\E_x[\cdot]$ for functions on
$\pmone^n$, we mean the expectation under the uniform measure: each
$x \in \pmone^n$ has probability $2^{-n}$. So
$\norm{f}_2^2 = \E_x[f(x)^2] = 2^{-n} \sum_{x \in \pmone^n} f(x)^2$.
\end{remark}


% ──────────────────────────────────────────────────────────────────────
\section*{Recommended Reading}

\begin{itemize}
  \item \textbf{\citet{axler2024linear}, \emph{Linear Algebra Done Right},
    4th ed.}
    A clean, determinant-free treatment of inner product spaces and the
    spectral theorem. Assumes only calculus. Chapters~6--7 cover exactly
    the inner-product-space machinery used in this guide.

  \item \textbf{\citet{horn2012matrix}, \emph{Matrix Analysis}, 2nd ed.}
    The standard reference for Kronecker products (Chapter~4), matrix
    norms, and structured matrices. More advanced than Axler; useful when
    you need precise eigenvalue bounds.

  \item \textbf{\citet{durrett2019probability}, \emph{Probability: Theory
    and Examples}, 5th ed.}
    A graduate-level probability text. Chapters~3--4 cover the CLT,
    Berry--Ess\'een, and convergence in distribution rigorously. Freely
    available from the author's website.

  \item \textbf{\citet{feller1971introduction}, \emph{An Introduction to
    Probability Theory and Its Applications}, Vol.~2.}
    The classic reference for characteristic-function proofs of the CLT
    and Berry--Ess\'een. Older style but very thorough. Good for seeing
    where the constant $C < 0.4748$ ultimately comes from.
\end{itemize}
